\section{Transfer Orbit Design}
\label{sec:transfer}

Transfer design in e2m2e covers impulsive, low-energy, and low-thrust
propulsion modes. The core methodology is a two-step ``search then optimize''
scheme: ballistic propagation on a coarse grid first filters the feasible
solution families, then the best candidates are handed to nonlinear
programming (NLP) for refinement. The unified design variables are the
tangential velocity ratio $\alpha$ (controlling the departure direction), the
transfer time $T$, and the time $t_{\mathrm{ins}}$ from apoapsis to the
insertion point, with the total objective
\begin{equation}
  J = \Delta v_1 + \Delta v_2 ,
\end{equation}
subject to position continuity at the insertion point, velocity-direction
parallelism with the target, and collision avoidance with the Earth and the
Moon. The dynamical baseline is the CR3BP, with $\mu$ taken from the DE421
constant baseline.

\subsection{Impulsive Transfers}

\subsubsection{Lambert Solving and Porkchop Scans}

The Lambert solver implements the Izzo (2015) algorithm with a Rust kernel
(Python only converts and wraps): it supports short-/long-way directions and
multi-revolution solutions (returning the low-energy right branch), and
provides a batch interface---a grid of $N$ geometries $\times\,M$ times of
flight enters Rust in one call, with infeasible combinations filled with NaN
without affecting the rest. The solver is verified against Vallado's classic
test case.

Solving Lambert point by point over a departure-epoch $\times$ time-of-flight
grid yields the porkchop two-impulse $\Delta V$ map; terminal states are
extracted through the terminal-condition abstraction
(\texttt{TerminalCondition}), and analytic terminals need no dynamics
propagation.

\subsubsection{Three-Body Targeting}

A two-body Lambert solution only qualifies as an initial guess in cislunar
space. The three-body Lambert targeter seeds a damped Newton iteration in the
CR3BP with the two-body solution: propagating with the state transition
matrix $\stm(t, t_0)$, each correction is solved from the $\stm_{rv}$ block
of the terminal STM; the step is halved whenever a full step increases the
error; convergence is declared when the terminal position error falls below
$10^{-8}$ (dimensionless).

\subsubsection{Multi-Impulse Optimization and the Lawden Primer-Vector Test}

For an $n$-impulse transfer between fixed endpoints, the decision variables
are only the intermediate nodes $[t_i, \vect{r}_i]$; arcs between adjacent
nodes are closed by Lambert solutions (optionally refined by three-body
targeting), and SLSQP minimizes the total $\Delta V$. Whether a two-impulse
solution can still be improved is answered by the Lawden primer-vector test:
from the transversality conditions $\vect{p}(t_0) = \hat{\Delta v}_0$ and
$\vect{p}(t_f) = \hat{\Delta v}_f$, the costate is carried by the STM to give
the $\vect{p}(t)$ curve, with the necessary optimality condition
\begin{equation}
  \norm{\vect{p}(t)} \le 1 \quad \text{throughout},
  \qquad \norm{\vect{p}} = 1 \text{ at impulses, aligned with them}.
\end{equation}
If $\norm{\vect{p}} > 1$ within an arc, inserting an intermediate impulse
reduces the total $\Delta V$ (Lion \& Handelsman, 1968); the test also
reports the suggested insertion time and position, directly seeding a
three-impulse optimization. A LEO$\to$GEO example reproduces the Hohmann
benchmark of $3.7708$ km/s.

The classical Hohmann two-impulse minimum-energy transfer
\begin{equation}
  \Delta v_1 = \sqrt{\frac{\mu}{r_1}}\left(\sqrt{\frac{2 r_2}{r_1+r_2}} - 1\right),
  \qquad
  \Delta v_2 = \sqrt{\frac{\mu}{r_2}}\left(1 - \sqrt{\frac{2 r_1}{r_1+r_2}}\right)
\end{equation}
is provided as an analytic benchmark and a source of initial guesses.

\subsection{Low-Energy Transfers}

\subsubsection{Lunar Gravity Assist (LGA)}

A three-leg conic-patching design: hyperbolic departure from the initial
orbit $\to$ a gravity assist inside the lunar sphere of influence that
rotates the velocity vector $\to$ hyperbolic insertion into the target orbit.
The search interface works in dimensionless CR3BP states; a typical search
domain is 15--45 days of flight time with a total $\Delta v$ cap of
4.0 km/s, and results report total cost, flight time, and perilune altitude.

\subsubsection{WSB Ballistic Capture}

The weak stability boundary (WSB) route exploits the Sun--Earth weakly stable
region: the spacecraft leaves the lunar sphere of influence onto a distant
Sun-perturbed arc, where solar gravity does work near apoapsis so that the
Moon-relative velocity drops below the capture threshold, yielding ballistic
capture at near-zero $\Delta v$ (Belbruno \& Miller, 1993). The capture
criterion is a negative two-body Kepler energy at perilune:
\begin{equation}
  H_2 = \frac{\norm{\vect{v}}^2}{2} - \frac{\mu_{\mathrm{moon}}}{r} < 0 .
\end{equation}
The solution proceeds in two steps: a \textbf{ballistic grid search}
propagates forward in the BCR4BP over a Sun-phase $\times$ departure-phase
$\times$ flight-time grid, filtering by perilune altitude and the $H_2$
criterion---propagation, section detection, and filtering all run in Rust
with Rayon parallelism (the Python reference implementation runs only when
explicitly requested, never as an automatic fallback); an \textbf{arrival-leg
refinement} then closes the best candidate's Moon-centered arrival arc with
three-body Lambert targeting. A typical search domain is 90--150 days.

\subsubsection{Invariant-Manifold Stitching}

A third low-energy route exploits the invariant manifold tubes of periodic
orbits themselves: crossings of the departure orbit's unstable manifold and
the target orbit's stable manifold on the same Poincar\'e section (a planar
section, or a periapsis section $s = \vect r \cdot \vect v = 0$) are paired
and ranked by the weighted patching cost
\begin{equation}
  w_r \norm{\Delta \vect r} + w_v \norm{\Delta \vect v}.
\end{equation}
The end-to-end pipeline: propagate the departure unstable manifold $\times$
target stable manifold (all four $\pm$ branch combinations, globally best
taken) to the secondary's periapsis section, pick the best patch, then close
onto the target orbit after the patch point with three-body Lambert
targeting; impulses accordingly split into departure, patch, and arrival
segments. In the benchmark case (medium-to-large-amplitude members of the L1
Lyapunov family), the patch impulse is on the order of tens of m/s. The route
is currently CR3BP-based, with an ephemeris interface reserved.

The value of low-energy routes: trading time for propellant, enabling
economical lander/orbiter scenarios and even emergency low-fuel rescue
corridors.

\subsection{Low-Thrust Transfers}

The low-thrust leg uses the 7-dimensional augmented state
$[\vect{r}, \vect{v}, m]$:
\begin{equation}
  \dot{\vect r} = \vect v, \qquad
  \dot{\vect v} = \vect a_{\mathrm{grav}} + \frac{T\,\delta}{m}\,\hat{\vect u}, \qquad
  \dot m = -\frac{T\,\delta}{I_{sp}\, g_0},
\end{equation}
with throttle $\delta \in [0,1]$ and thrust direction $\hat{\vect u}$
parameterized by spherical angles $(\theta_1, \theta_2)$. Solving is a
two-level loop: a \textbf{Q-law initial guess} integrates a Lyapunov feedback
law forward to produce a suboptimal control history; an \textbf{SLSQP
polish} then minimizes propellant over piecewise-constant controls
$(\delta, \theta_1, \theta_2)$ per segment, matching the terminal position
and velocity. Two solvers are available: shooting (analytic Jacobians,
5--24$\times$ faster) and Hermite--Simpson collocation (more robust at large
scale).

Impulsive and low-thrust costs are distinguished via the equivalent
$\Delta v$: impulsive uses $\Delta v = \sum_i \norm{\Delta \vect v_i}$, while
low-thrust inverts the Tsiolkovsky equation,
$\Delta v_{\mathrm{eq}} = I_{sp}\, g_0 \ln(m_0/m_f)$; due to gravity losses,
$\Delta v_{\mathrm{eq}} \ge$ the impulsive $\Delta V$, with significant
differences on high-energy transfers.

\subsection{The Grid-Search + NLP Two-Step Optimization Framework}

Step one, \textbf{grid search}: departure points are sampled equally in time
along the departure orbit (typically 200), each with $\alpha$ uniformly
sampled in $[0.5, 2.5]$ (typically 101); transfer arcs are integrated forward
and filtered by collision (200 km at Earth / 100 km at the Moon),
intersection, and minimum-distance ($10^{-3}$) constraints; the numerical
kernel has sunk to Rust and runs in parallel.

Step two, \textbf{NLP refinement}: decision variables
$\vect{y} = \{\alpha, T, t_{\mathrm{ins}}\}$, objective
$J(\vect y) = \Delta v_1 + \Delta v_2$, constraints of insertion-point
position continuity (equality) and velocity parallelism (equality, or relaxed
to an inequality with an angular tolerance of 0.05 rad). The solver defaults
to SciPy SLSQP, with the COPT commercial optimizer as an option and fallback
on failure.

\begin{table}[htbp]
\centering\small
\caption{Transfer design capabilities at a glance}
\label{tab:transfer}
\begin{tabularx}{\textwidth}{@{}llX@{}}
\toprule
\textbf{Category} & \textbf{Method} & \textbf{Notes} \\
\midrule
Impulsive & Lambert (Izzo) & Rust kernel, batch grid solving, short/long
way and multi-revolution \\
Impulsive & Three-body targeting & Damped Newton + STM, position error
$< 10^{-8}$ at convergence \\
Impulsive & Multi-impulse & SLSQP + Lawden primer-vector test, automatic
impulse-insertion suggestions \\
Impulsive & Hohmann & Analytic benchmark and initial guesses \\
Low-energy & LGA & Three-leg conic patching, TOF 15--45 days \\
Low-energy & WSB & BCR4BP + $H_2<0$ criterion, Rust Rayon parallel search \\
Low-energy & Manifold stitching & Poincar\'e section pairing + three-body
closure, patch impulses at m/s scale \\
Low-thrust & Q-law + NLP & Shooting (analytic Jacobian, 5--24$\times$) or
collocation \\
Framework & Search--optimize & Grid search (Rust parallel) $\to$ SLSQP/COPT
refinement \\
\bottomrule
\end{tabularx}
\end{table}
